% Generated draft: requires mathematical and source review.
\documentclass[11pt]{article}
\usepackage[margin=25mm]{geometry}
\usepackage{fontspec}
\setmainfont{Latin Modern Roman}
\usepackage{amsmath,amssymb,mathtools,unicode-math}
\setmathfont{Latin Modern Math}
\usepackage{enumitem,needspace,xurl,hyperref}
\setlength{\parindent}{0pt}
\setlength{\parskip}{5pt}
\setlength{\emergencystretch}{3em}
\newcommand{\N}{\mathbb{N}}
\newcommand{\Z}{\mathbb{Z}}
\newcommand{\Q}{\mathbb{Q}}
\newcommand{\R}{\mathbb{R}}
\newcommand{\C}{\mathbb{C}}
\newcommand{\F}{\mathbb{F}}
\newcommand{\A}{\mathbb{A}}
\newcommand{\PP}{\mathbb{P}}
\newcommand{\E}{\mathbb{E}}
\newcommand{\eps}{\varepsilon}
\newcommand{\dd}{\,\mathrm{d}}
\begin{document}
\section*{UnsolvedMath ID: 2630}
\textbf{Source problem code:} catalogue entry
\subsection{Definitions and mathematical statement}
Fix a prime $p$. For a finite group $G$, let $G(p)$ mean any Sylow $p$-subgroup, so $|G(p)|$ is the largest power of $p$ dividing $|G|$, independent of the chosen Sylow subgroup. This interpretation makes explicit notation left undefined in the source.

Using nonnegative real exponents, define the admissible-exponent set
\[
 \mathcal E_p(\Gamma)=\left\{e\in[0,\infty):\ \exists J\ge1\
 \forall G\le\Gamma\text{ finite}\ \exists A\triangleleft G,
 \ \begin{gathered}
 A\text{ abelian},\quad p\nmid|A|,\\
 [G:A]\le J|G(p)|^e
 \end{gathered}\right\}.
\]
The group $\Gamma$ is $p$-Jordan if this set is nonempty, and its exponent is
\[
 e_p(\Gamma)=\inf\mathcal E_p(\Gamma).
\]
For fixed $p,\Gamma,e$, the finite constant $J$ may depend on all three but must work for every finite subgroup $G$. The normal abelian subgroup $A$ may depend on $G$, and the trivial subgroup is permitted. No linearity, finite-generation or topological assumption is imposed on $\Gamma$.

The two questions are:
\begin{enumerate}
\item For every $p$-Jordan $\Gamma$, is $e_p(\Gamma)\in\mathcal E_p(\Gamma)$? Equivalently, does one finite constant $J$ make the bound valid at the infimum itself?
\item For every $p$-Jordan $\Gamma$, is $e_p(\Gamma)\le3$?
\end{enumerate}
The second, as written, is an infimum bound. Without the first, it does not necessarily assert that exponent exactly $3$ is admissible in the boundary case $e_p(\Gamma)=3$. The admissible set is upward closed, so an infimum at most $3$ gives all strictly larger exponents, but endpoint attainment is a separate issue.

The source does not explicitly specify the allowed domain of $e$. The displayed definition uses the customary nonnegative real domain needed for this meaningful attainment question. If arbitrary negative real exponents were allowed, the trivial group would admit every real exponent and have infimum $-\infty$, providing a degenerate nonattainment issue. If only nonnegative integers were allowed, the infimum question would instead be automatically attained by well-ordering. The real nonnegative convention therefore requires source confirmation, rather than being treated as inconsequential.

The source reports that for every field $K$ of characteristic $p$ and every fixed $n\ge1$, results of Brauer--Feit and Larsen--Pink give $3\in\mathcal E_p(\operatorname{GL}_n(K))$. This is an admissible exponent bound, not a claim that the optimal infimum is always exactly $3$. The general questions concern all $p$-Jordan groups, beyond this linear benchmark. The literature triage checked 17 August 2026 says both attainment and the universal infimum bound remain open, classifies the record partially solved by the linear examples, and attributes the question to C. Shramov in Kourovka Notebook 21.121. The main registry status is open.
\subsection{Short English statement}
For a group whose finite subgroups admit normal abelian prime-to-$p$ subgroups with a controlled index, is the optimal real exponent attained, and is its infimum always at most $3$? The exponent domain and Sylow notation need explicit conventions; attainment at $3$ is separate from an infimum bound.
\subsection{Sources}
\begin{itemize}
\item Local UnsolvedMath registry, permanent ID $2630$, exact source code \texttt{catalogue entry}, ``Kourovka Notebook Problem 21.121,'' including the supplied literature triage checked 17 August 2026.
\item \emph{Kourovka Notebook}, New Problems, 21st issue (2026), Problem 21.121, attributed to C. Shramov, \url{https://www.kourovka-notebook.org/}.
\item The statement attributes the exponent-$3$ linear-group benchmark to Brauer--Feit and Larsen--Pink, but supplies no individual paper titles, publication details or URLs; no missing references are invented.
\end{itemize}
\end{document}
